\section{Discussion}
\label{sec:cls_discussion}
Bone-suppressed inputs did not improve these classifiers. The multi-stream hybrid matched the radiograph alone on all three datasets, and on
NIH the interval rules out a gain larger than 0.0024 in macro-AUROC, so the null result is informative rather than a lack of power. The
hybrid also matched the control that gives the same network the radiograph three times, and plain channel stacking did at least as well as
separate streams. With 175 runs, seed ensembles and a plan fixed before testing, these results leave little room for a general accuracy
gain from suppression with this suppressor and encoder.

The information argument (Section~\ref{sec:cls_intro}) offers one explanation. The suppressor adds no information, only a transformation learned from
CT, and a DenseNet-121 trained on thousands of radiographs may learn whatever bone-invariant features its labels reward. The prior
can matter only where the classifier cannot easily learn the separation itself, and the finding groups suggest where that is.

\paragraph{Focal lesions.} Nodules and masses had the largest estimates on both adult datasets, of similar size ($+0.008$ to $+0.009$ for
the soft-tissue input), and were the only group with a gain on NIH. Per class, the NIH gain came from nodules; masses did not change. This
fits reader studies in which suppression improved nodule detection \citep{freedman2011lung}, and the geometry: a rib can hide a small
nodule more easily than a mass. The pointing game agrees: maps of the radiograph stream found focal lesions less often than other findings,
while maps of the soft-tissue stream found both about equally often, a pattern that Integrated Gradients reproduced. Two cautions limit this result. The NIH test split was used in \companion{ch:benchmark}, whose
classifiers already showed a nodule gain and losses for effusion, infiltration and pleural thickening on the same images; the NIH result
is therefore a re-test with new models rather than independent confirmation. VinBigData, the independent test, has a single merged
nodule/mass class, and its interval included zero.

\paragraph{Diffuse findings.} The loss for diffuse findings appeared on NIH only and came mostly from infiltration. NIH labels were mined
from reports \citep{wang2017chestxray}, and infiltration is a broad label; on VinBigData, labelled by radiologists,
diffuse findings did not change. These data cannot tell whether suppression removes image cues for diffuse disease or interacts with label
noise.

\paragraph{Uncertainty.} The conformal procedures kept their guarantees on all inputs, but NIH label sets that hold half of the 14 labels
are of little use to a reader. The suppressor's own uncertainty made abstention worse in every configuration. Uncertainty about the bone image is a different quantity from uncertainty about the diagnosis, and this result argues
against using one for the other.

\paragraph{Implications.} A classifier should not be expected to gain overall accuracy from bone-suppressed inputs. Where suppression is
considered for focal lesions, it should be evaluated per finding type and against a capacity-matched control, since a multi-stream network
costs three times the computation of a single-stream one. Suppression for human readers is a separate question that these results do not
address.

\subsection{Limitations}
\label{sec:cls_limits}
\begin{itemize}
\item \emph{One encoder.} The pre-specified screen selected DenseNet-121 alone; heavier encoders may use the extra streams differently.
\item \emph{Resolution.} Classifiers saw $256\times256$\,px images, where small nodules occupy few pixels; suppression may matter more at
  full resolution.
\item \emph{One suppressor.} The suppressor was trained to keep lesion signal in the soft-tissue image; a suppressor that absorbs lesions
  could remove the focal gain or turn it into a loss.
\item \emph{Reused test split and labels.} The NIH test split had been used with other models in \companion{ch:benchmark}, and its labels are
  report-derived. VinBigData has one focal class, and BDCXR-3257 is a three-class pneumonia task without focal findings.
\item \emph{Clustering in per-class tests.} The pre-specified DeLong tests treat NIH images as independent although patients contribute
  several images; the group and overall tests use a patient-cluster bootstrap.
\item \emph{Attribution.} The pointing game depends on the attribution method and the 15\,px tolerance, and the expert boxes cover eight
  NIH classes only.
\item \emph{Registration.} The plan was hash-frozen before testing but not deposited in a public registry.
\end{itemize}

\section{Conclusion}
Bone-suppressed inputs, alone, stacked or as separate streams, did not make deep chest radiograph classifiers more accurate on three
datasets. The largest effect on both adult datasets was a small gain for nodules and masses, significant on NIH only and matched by better
localisation in the soft-tissue stream; on NIH it came with a loss for diffuse findings. Bone suppression is not a general accuracy gain for classifiers, and
any benefit should be sought, and tested, for focal lesions specifically.
